Beim Verbrennen wird pro Kilogramm Butan die Wärme $H = \HuO$ frei (spezifischer Heizwert, FoTaBe S.~193), insgesamt also $Q = m\,H$. Gibt der Brenner diese Wärme mit der Leistung $P$ ab, so reicht die Kartusche für
\begin{align}
 t &= \frac{Q}{P} = \tBrF \\
   &= \frac{\mBu\cdot\Hu}{\PBr} \\
   &= \tBr = \tBrh \approx \result{\tBrhP{0}{2}}
\end{align}
